\documentclass[11pt]{article}
\usepackage[a4paper,margin=24mm]{geometry}
\usepackage[T1]{fontenc}
\usepackage{lmodern,amsmath,amssymb,amsthm,mathtools,microtype,booktabs}
\usepackage[hidelinks]{hyperref}
\usepackage{xurl}
\hypersetup{pdftitle={Moment domination and a certified half-radius lower bound for the signed prime--gamma form},pdfauthor={Jongmin Choi}}
\newtheorem{theorem}{Theorem}
\newtheorem{corollary}[theorem]{Corollary}
\newtheorem{lemma}[theorem]{Lemma}
\newcommand{\R}{\mathbb R}
\newcommand{\dd}{\,\mathrm d}
\newcommand{\norm}[1]{\lVert#1\rVert}
\newcommand{\Pole}{\operatorname{Pole}}
\DeclareMathOperator{\supp}{supp}
\DeclareMathOperator{\Var}{Var}
\setlength{\emergencystretch}{2em}
\allowdisplaybreaks[1]
\title{Moment domination and a certified half-radius lower bound for the signed prime--gamma form}
\author{Jongmin Choi\\
Independent Researcher, Seoul, Korea\\
\texttt{24ping@naver.com}\\
ORCID iD: 0009-0008-7448-514X}
\date{27 October 2026 (UTC)\quad Version 0.3}
\begin{document}
\maketitle
\begin{abstract}
We prove explicit sufficient domination inequalities for a signed
prime--gamma quadratic form on compact physical supports. A finite set
of exponentially weighted cell moments controls the low-frequency
resolvent norm, yielding both nonlinear and quadratic lower bounds.
The complete gamma energy, every active prime-power correlation and
the pole are retained. The moment kernel has an explicit positive gap
for sufficiently fine cells. A finite sufficient matrix condition
retains its coupling with the remaining directions. A computer-assisted Schur estimate additionally proves the full-domain
bound $q_B(u)\ge7\cdot10^{-7}\lVert u\rVert_2^2$ for $B\le1/2$,
including the active prime $2$, the full gamma energy and pole, and an
analytic bound for every omitted cosine mode. Its finite calculation
uses validated ball arithmetic. A separate conditional application
assumes a specified localization inequality;
that assumption is not proved here. Positivity at arbitrary radii, a global
spectral lower bound $1/2$, and the Riemann hypothesis are not established.
No historical novelty claim is made.
\end{abstract}

\section{The signed form and its domain}
For real even compact smooth $u$, define
\begin{align}
 \rho(s)&=\frac{e^{-s/2}}{1-e^{-2s}},\qquad
 a_n=\frac{\Lambda(n)}{\sqrt n},\qquad C(x)=\cosh(x/2),\nonumber\\
 t_\Gamma(u)&=\int_0^\infty\rho(s)\norm{u(\cdot+s)-u}_2^2\dd s,\nonumber\\
 p_\infty(u)&=2\sum_{n\ge2}a_n\int_\R u(x)u(x+\log n)\dd x,\nonumber\\
 c_\Gamma&=\log\pi+\gamma_E+\pi/2+3\log2,\qquad
 \Pole(u)=2\left|\int_\R Cu\right|^2,\nonumber\\
 q_\infty(u)&=t_\Gamma(u)-c_\Gamma\norm{u}_2^2-p_\infty(u)+\Pole(u).
 \label{eq:form}
\end{align}
Here $\gamma_E$ is Euler's constant and $\Lambda$ is the von Mangoldt
function: $\Lambda(p^k)=\log p$ for prime $p$ and integer $k\ge1$,
and $\Lambda(n)=0$ otherwise. Positive gamma shifts are counted once.
The prime factor two counts both orientations. All prime powers retain
their actual von Mangoldt weights.

Fix $B>0$ and work with the minimum compact-support form domain $\mathcal D_B$, the closure of real even $C_c^\infty((-B,B))$, zero extended, in $\bigl(\norm{u}_2^2+t_\Gamma(u)\bigr)^{1/2}$. The gamma form is closable, as follows from its nonnegative Fourier multiplier below. Every prime correlation with $\log n\ge2B$ is exactly zero, including equality. This is support geometry, not a discarded prime tail. The finite prime operator and pole are bounded on this compact $L^2$ space; consequently the resulting semibounded form is closed on $\mathcal D_B$, by the standard closed-form framework \cite[Chapter VI]{Kato}.

Set
\begin{equation}
 W_r^{\rm p}=\sum_{2\le n\le e^{2r}}a_n,\qquad
 \beta_B=W_B^{\rm p}+W_{B/2}^{\rm p}.
 \label{eq:beta}
\end{equation}
The letter $W_r^{\rm p}$ denotes a finite prime-weight sum. It is distinct from the auxiliary weight introduced in Section~\ref{sec:localization}.

\section{The finite moments and the inequality}
Partition $(0,B)$ into $m$ intervals of positive lengths at most $\delta$, and reflect them to $(-B,0)$. In the real even ordinary-$L^2$ space let $Z_\delta$ be the span of the Riesz representers of all cell functionals
\[
 u\longmapsto\int_J e^{x/2}u(x)\dd x.
\]
Equivalently, each positive cell contributes the two even extensions of $e^{x/2}$ and $e^{-x/2}$ on that cell. Thus $\dim Z_\delta=2m$. Let $P_\delta$ be the ordinary-$L^2$ orthogonal projection onto $Z_\delta$, and define
\begin{equation}
 \eta_\delta=\frac1{1+(2\pi/\delta)^2},\qquad
 \sigma(u)=\frac{\norm{P_\delta u}_2}{\norm u_2}\in[0,1]\quad(u\ne0),
 \label{eq:eta}
\end{equation}
\begin{equation}
 s_\delta(u)=\min\left\{1,
 \left[\sigma(u)+\sqrt{\eta_\delta}\sqrt{1-\sigma(u)^2}\right]^2\right\}>0.
 \label{eq:s}
\end{equation}

\begin{theorem}[Moment-defect domination]
Every nonzero $u\in\mathcal D_B$ satisfies
\begin{equation}
 \boxed{q_\infty(u)\ge
 \left[-\tfrac12\log s_\delta(u)-c_\Gamma-\beta_B\right]\norm u_2^2
 +\Pole(u).}
 \label{eq:nonlinear}
\end{equation}
There is also the simpler quadratic bound, valid including $u=0$:
\begin{equation}
 \boxed{q_\infty(u)\ge A_{B,\delta}\norm u_2^2
 -K_\delta\norm{P_\delta u}_2^2+\Pole(u),}
 \label{eq:linear}
\end{equation}
where
\begin{equation}
 A_{B,\delta}=\frac12\log\frac{1+(2\pi/\delta)^2}{2}-c_\Gamma-\beta_B,
 \qquad K_\delta=\frac{2\pi^2}{\delta^2}.
 \label{eq:constants}
\end{equation}
\end{theorem}

\begin{corollary}[Explicit winning conditions]
If $g>0$ and
\begin{equation}
 \sigma(u)+\sqrt{\eta_\delta}\sqrt{1-\sigma(u)^2}
 \le e^{-(c_\Gamma+\beta_B+g)},
 \label{eq:win}
\end{equation}
then $q_\infty(u)\ge g\norm u_2^2+\Pole(u)$. Alternatively, if $A_{B,\delta}>0$ and
\[
 \norm{P_\delta u}_2^2\le\frac{A_{B,\delta}}{2K_\delta}\norm u_2^2,
\]
then $q_\infty(u)\ge\tfrac12A_{B,\delta}\norm u_2^2+\Pole(u)$.

If the cell moments vanish exactly, the sharper bound is
\begin{equation}
 q_\infty(u)\ge
 \left[\tfrac12\log\bigl(1+(2\pi/\delta)^2\bigr)-c_\Gamma-\beta_B\right]\norm u_2^2.
 \label{eq:kernel}
\end{equation}
Here $\Pole(u)=0$ follows from the same moments. For any prescribed $g>0$, the coefficient in \eqref{eq:kernel} is at least $g$ whenever
\begin{equation}
 \delta\le\frac{2\pi}{\sqrt{e^{2(c_\Gamma+\beta_B+g)}-1}}.
 \label{eq:deltachoice}
\end{equation}
This gives an explicit infinite-dimensional class on which the positive gamma energy dominates the entire negative mass-plus-prime expression.
\end{corollary}

\section{Proof, including cross terms}
Use $\widehat u(\xi)=\int u(x)e^{-i\xi x}\dd x$. Define the bounded positive form
\begin{equation}
 J(u)=\frac1{2\pi}\int_\R\frac{|\widehat u(\xi)|^2}{1+4\xi^2}\dd\xi
 =\frac14\iint e^{-|x-y|/2}u(x)u(y)\dd x\dd y.
 \label{eq:J}
\end{equation}
It obeys $0<J(u)\le\norm u_2^2$ for $u\ne0$.

\paragraph{Step 1: the gamma lower bound.}
With $a_j=2j+1/2$, expansion of the positive geometric series gives
\[
 d_\Gamma(\xi)=\sum_{j\ge0}\frac{2\xi^2}{a_j(a_j^2+\xi^2)}
 \ge\frac12\log(1+4\xi^2).
\]
For $\xi\ne0$ the summand, viewed as a function of $a>0$, decreases. Compare the sum with the integral on $a\ge1/2$ with mesh factor $1/2$; the integral is exactly the displayed logarithm. At zero both sides vanish. Also
\[
 d_\Gamma(\xi)\le 2\xi^2\sum_{j\ge0}(2j+1/2)^{-3}
 \le\left(16+\tfrac14\sum_{j\ge1}j^{-3}\right)\xi^2
 \le\tfrac{131}{8}\xi^2<18\xi^2\quad(\xi\ne0),
\]
where $\sum_{j\ge1}j^{-3}\le1+\int_1^\infty t^{-3}\dd t=3/2$.
Plancherel and Jensen's inequality for $-\log$ (see \cite{HLP}), with probability density $|\widehat u|^2/(2\pi\norm u_2^2)$, imply
\begin{equation}
 t_\Gamma(u)\ge\frac{\norm u_2^2}{2}
 \log\frac{\norm u_2^2}{J(u)}.
 \label{eq:Jensen}
\end{equation}
Finite gamma energy and a standard truncation justify Jensen. This is a lower bound on the complete gamma energy, not on a finite shift truncation.

\paragraph{Step 2: why vanishing moments force coercivity.}
Write $u=v+w$, with $w=P_\delta u$ and $v=(I-P_\delta)u$. These functions are ordinary-$L^2$ orthogonal. Every cell moment of $v$ is zero. On $J=(a,b)$ set
\[
 k(x)=e^{-x/2}\int_a^x e^{y/2}v(y)\dd y.
\]
Then $k(a)=k(b)=0$ and $v=k'+k/2$. Joining the zero extensions gives $k\in H^1(\R)$. The cross term integrates to zero on each cell. The one-dimensional Dirichlet Poincare inequality, equivalently the endpoint-zero Wirtinger inequality \cite{HLP}, gives
\[
 \norm v_2^2=\norm{k'}_2^2+\tfrac14\norm k_2^2
 \ge\bigl[(\pi/\delta)^2+1/4\bigr]\norm k_2^2.
\]
For completeness, on a cell of length $\ell$ the sine expansion of an $H_0^1$ function has derivative energy $\sum_{n\ge1}(n\pi/\ell)^2|c_n|^2$ and squared norm $\sum_{n\ge1}|c_n|^2$, proving the constant $(\pi/\ell)^2\ge(\pi/\delta)^2$. Fourier transformation of $v=k'+k/2$ yields $J(v)=\norm k_2^2/4$, hence
\begin{equation}
 J(v)\le\eta_\delta\norm v_2^2.
 \label{eq:kernelJ}
\end{equation}
No regularity of the orthogonal projection in the unbounded gamma operator domain is needed: this step concerns only $L^2$ functions and the bounded form $J$.

\paragraph{Step 3: retain the mixed term.}
The norm $J^{1/2}$ is a Hilbert norm, so
\[
 \sqrt{J(u)}\le\sqrt{J(v)}+\sqrt{J(w)}
 \le\sqrt{\eta_\delta}\norm v_2+\norm w_2.
\]
Using $\norm v_2^2+\norm w_2^2=\norm u_2^2$, and also $J(u)\le\norm u_2^2$, proves
$J(u)/\norm u_2^2\le s_\delta(u)$. Substitution into \eqref{eq:Jensen} gives the nonlinear gamma part of \eqref{eq:nonlinear}.

For the quadratic form version, the elementary squared triangle bound gives
\[
 J(u)\le2\eta_\delta\norm u_2^2
             +2(1-\eta_\delta)\norm{P_\delta u}_2^2.
\]
For any $z,a>0$, $-\log z\ge-\log a+1-z/a$. Set
$z=J(u)/\norm u_2^2$ and $a=2\eta_\delta$ in \eqref{eq:Jensen}. This gives
\[
 t_\Gamma(u)\ge\frac12\log\frac1{2\eta_\delta}\norm u_2^2
 -\frac{1-\eta_\delta}{2\eta_\delta}\norm{P_\delta u}_2^2.
\]
The last coefficient is exactly $K_\delta=2\pi^2/\delta^2$. In particular, the mixed term has been bounded, not dropped.

\paragraph{Step 4: every active prime power.}
On $(-B,B)$ let
\[
 (\mathsf P_Bu)(x)=\sum_{\log n<2B}a_n
       \{u(x+\log n)+u(x-\log n)\},
\]
using zero extensions. This is a symmetric finite positive-kernel operator and $p_\infty(u)=\langle u,\mathsf P_Bu\rangle_2$. For $x\ge0$, the two available travel lengths are at most $B+x\le2B$ and $B-x\le B$; for $x<0$ reflect. Thus every row sum is at most $\beta_B$. Applying $2|u(x)u(y)|\le |u(x)|^2+|u(y)|^2$ to each shift and using symmetry bounds the absolute quadratic form by the integral of the row sum times $|u(x)|^2$. This proves the Schur-type estimate
\begin{equation}
 |p_\infty(u)|\le\beta_B\norm u_2^2.
 \label{eq:primebound}
\end{equation}
Adding the unchanged local mass and pole to the gamma lower bounds proves both inequalities. All bounded terms pass to $\mathcal D_B$ by continuity. The gamma form passes by its defining form convergence. The winning conditions follow by substitution.

\paragraph{Nonvacuity.}
The moment kernel is infinite-dimensional. For example, on any one positive cell choose $\phi\in C_c^\infty(J)$ and put $v=(\partial_x^2-1/4)\phi$ there, extending evenly to the negative cell and by zero elsewhere. Integration by parts shows that both $e^{x/2}$ and $e^{-x/2}$ moments vanish. Such nonzero functions form an infinite-dimensional family. This is an exact construction, not a numerical approximate nullspace.

\section{An explicit finite-moment observable}
For a positive cell $J=(a,b)$, let the two basis functions be the even extensions of $e^{x/2}\mathbf1_J$ and $e^{-x/2}\mathbf1_J$. Their ordinary-$L^2$ Gram matrix and moment vector are
\[
 G_J=2\begin{pmatrix}e^b-e^a&b-a\\ b-a&e^{-a}-e^{-b}\end{pmatrix},
 \qquad z_J=2\begin{pmatrix}\int_J e^{x/2}u(x)\dd x\\
                              \int_J e^{-x/2}u(x)\dd x\end{pmatrix}.
\]
Then
\begin{equation}
 \norm{P_\delta u}_2^2=\sum_{J\subset(0,B)}z_J^{\mathsf T}G_J^{-1}z_J.
 \label{eq:observable}
\end{equation}
For $\ell=b-a>0$, $\det G_J=4(2\cosh\ell-2-\ell^2)>0$, by the positive Taylor terms after degree two. Thus the winning condition uses an explicitly defined finite set of moments and $2\times2$ matrices, not an unknown signed-form minimum. Direct evaluation on tiny cells should use stable hyperbolic expressions or validated arithmetic; no floating evaluation is needed for this theorem.

\section{A conditional localization application}\label{sec:localization}
Let $W$ be a continuous positive even weight, and put
$N_X(u)=\int_\R |u|^2/W$. On compact intervals this norm and the ordinary
$L^2$ norm are equivalent. For $R\ge1$ define
\[
 b(t)=\begin{cases}0,&t\le0,\\3t^2-2t^3,&0<t<1,\\1,&t\ge1,\end{cases}
 \qquad \theta_R(x)=\tfrac\pi2b(|x|-R),
 \quad\chi_R=\cos\theta_R,\quad\eta_R=\sin\theta_R.
\]
These even $C^1$ cutoffs satisfy $\chi_R^2+\eta_R^2=1$,
$\chi_R=1$ on $[-R,R]$ and $\chi_R=0$ outside $[-R-1,R+1]$.
The cutoff $\eta_R$ is distinct from the scalar $\eta_\delta$.

\paragraph{Localization hypothesis.}
Assume that for every $R\ge1$ there is an $\varepsilon_R\ge0$ such that,
on the real even compact smooth core,
\begin{equation}
 q_\infty(u)\ge q_\infty(\chi_Ru)+t_\Gamma(\eta_Ru)
                   -\varepsilon_RN_X(u).
 \label{eq:localization-hypothesis}
\end{equation}
This is an explicit additional assumption. No bound for
$\varepsilon_R$, or existence of such a bound for a particular weight,
is proved in this paper. In particular, the elementary partition identity
does not imply \eqref{eq:localization-hypothesis} by itself. Any proof of
this hypothesis must control the complete mass, prime and pole defects.

Under this hypothesis, choose $B=R+1$ and the corresponding moment
projection. The cutoff product belongs to $\mathcal D_B$: it is compact
$H^1$, vanishes at the endpoints, and is approximable by interior smooth
tests in $H^1$. The bound $d_\Gamma(\xi)\le18\xi^2$ makes that
approximation converge in the gamma-form norm. Substitution of
\eqref{eq:linear} into \eqref{eq:localization-hypothesis} gives
\begin{equation}
 \boxed{\begin{aligned}
 q_\infty(u)\ge{}&A_{B,\delta}\norm{\chi_Ru}_2^2+t_\Gamma(\eta_Ru)
                  +\Pole(\chi_Ru)\\
                &-K_\delta\norm{P_\delta(\chi_Ru)}_2^2
                  -\varepsilon_RN_X(u).
 \end{aligned}}
 \label{eq:global}
\end{equation}
This conditional conclusion is a sufficient comparison, not an assertion
that its right side is nonnegative for every input. The sharper nonlinear
bound can replace its first and fourth terms. No prime tail is omitted
inside $[-B,B]$.

Even when $\int Cu=0$, the localized mean need not vanish. The possibly
negative difference $\Pole(u)-\Pole(\chi_Ru)$ must be covered by
\eqref{eq:localization-hypothesis}. The positive local pole in
\eqref{eq:global} is the one inside $q_\infty(\chi_Ru)$, not a second
use of the original pole.

\section{A finite sufficient test with the coupling retained}
The kernel bound supplies a concrete next proof path. Fix $B,\delta$ so that
\eqref{eq:kernel} gives $q_B(h)\ge g\norm h_2^2$ on $\ker P_\delta$,
with $g>0$ and $q_B$ the restriction of $q_\infty$ to $\mathcal D_B$.
Choose an ordinary-$L^2$ orthonormal basis $e_1,\ldots,e_d$ of $Z_\delta$
and write $L_0u=(\langle e_j,u\rangle_2)_{j=1}^d$.
Choose smooth even compactly supported lift columns $V_j$ satisfying
$L_0V_j$ equal to the $j$th coordinate vector. Such lifts exist: the cell
moment functionals are linearly independent on the smooth core, by its
density in the compact even $L^2$ space.

Let $L_B$ be the self-adjoint operator of the closed semibounded form
$q_B$, supplied by the representation theorem \cite[Chapter VI]{Kato}.
Smooth compactly supported columns are in its operator domain:
the gamma multiplier grows at most quadratically and the remaining terms
are bounded on this compact interval. Define
\[
 G_{ij}=q_B(V_i,V_j),\qquad
 r_j=(I-P_\delta)L_BV_j,\qquad
 H_{ij}=\langle r_i,r_j\rangle_2.
\]
For $u=Vz+h$, $z=L_0u$ and $h\in\ker P_\delta$, completion of the square
gives the explicit inequality
\begin{equation}
 \boxed{q_B(u)\ge
 z^{\mathsf T}(G-g^{-1}H)z
 +g\left\|h+g^{-1}\sum_jz_jr_j\right\|_2^2.}
 \label{eq:finite-sufficient}
\end{equation}
In particular, the finite matrix condition $G-g^{-1}H\succeq0$ is
sufficient for positivity of the full infinite-dimensional local form.
This cell-moment matrix condition is not established here.
Section~\ref{sec:half} certifies a different, frequency-dependent
sufficient Schur condition for $B=1/2$. The square is an $L^2$ square;
its individual summands need not lie in the form domain. The negative
matrix $-g^{-1}H$ is precisely the allowed cost of the mixed term.
Checking $G\succeq0$ alone does not suffice.

More precisely, let $B_K$ be the positive self-adjoint operator of the
restricted form on $K=\ker P_\delta$. The exact Schur complement is
$G-R^*B_K^{-1}R$, where $Rz=\sum_jz_jr_j$. It is bounded below by
$G-g^{-1}H$ since $B_K\ge gI$. Thus \eqref{eq:finite-sufficient} is a
conservative sufficient test, not an identification with the exact
Schur complement. The entries, residual norms and the lower bound $g$
must all refer to the same form and the same ordinary-$L^2$ metric.

For a possible localization limit, perform the same construction for
$q_B+\delta_B N_X$ with $\delta_B\ge0$. On fixed compact supports the
added weight is a bounded positive multiplication operator. A verified
positive sufficient matrix would prove
$q_B\ge-\delta_BN_X$ on the entire local domain. Obtaining such matrices
for an unbounded sequence of radii with $\delta_B\to0$ remains open.
As proved below, the resulting family of all-domain inequalities is
equivalent to positivity on every compact local domain, rather than
a weaker intermediate objective.

\section{A certified lower bound on the whole half-radius domain}
\label{sec:half}
The next result treats the full local domain, including nonzero cell
moments. Its finite calculation uses validated ball arithmetic. It does
not assume the localization hypothesis.

\begin{theorem}[Computer-assisted half-radius coercivity]\label{thm:half}
For the normalization in \eqref{eq:form}, every real even
$u\in\mathcal D_B$, $0<B\le1/2$, satisfies
\begin{equation}\label{eq:halfwin}
 \boxed{q_B(u)\ge\frac7{10^7}\norm u_2^2.}
\end{equation}
The analytic reduction and tail estimates are proved below. The finite
matrix inequalities are certified by the accompanying
\texttt{verify\_half\_support.py}, with its recorded interval output.
This is a computer-assisted result with an explicit software trust
boundary; it is not a proof-assistant verification or independent review.
\end{theorem}

\subsection{Exact Neumann representation of the gamma energy}
For now take any $B>0$ and set $f(x)=\sqrt2u(x)$ on $(0,B)$.
This is a unitary identification of the even ordinary-$L^2$ space.
Let
\[
 k_n=\frac{n\pi}{B},\quad b_n=\sqrt{\frac{2-\delta_{n0}}B},\quad
 e_n(x)=b_n\cos(k_nx),\quad f_n=\langle f,e_n\rangle.
\]
Write $\alpha_j=2j+1/2$; this notation separates gamma labels from
prime weights. The complete gamma symbol is
\begin{equation}\label{eq:exact-symbol}
 d_\Gamma(k)=\sum_{j\ge0}\frac{2k^2}{\alpha_j(\alpha_j^2+k^2)}
 =\Re\psi(1/4+ik/2)-\psi(1/4).
\end{equation}
Here $\psi=\Gamma'/\Gamma$. The identity follows directly from its
convergent partial-fraction series \cite[\S5.7(ii)]{DLMF}.
Using $\psi(1/4)=-\gamma_E-\pi/2-3\log2$, define
\[
 D_n=d_\Gamma(k_n)-c_\Gamma
     =\Re\psi(1/4+ik_n/2)-\log\pi.
\]

For $a>0$, let $A_N=-\partial_x^2$ on $(0,B)$ with Neumann endpoints.
Its resolvent kernel gives the exact identity
\begin{equation}\label{eq:kernel-identity}
 e^{-a|x-y|}+e^{-a(x+y)}
 =2a(A_N+a^2)^{-1}(x,y)
   -2(\coth(aB)-1)\cosh(ax)\cosh(ay).
\end{equation}
Indeed for $x\le y$ the resolvent kernel times $2a$ is
$2\cosh(ax)\cosh(a(B-y))/\sinh(aB)$; subtraction verifies the formula.
Since $\rho(s)=\sum_{j\ge0}e^{-\alpha_js}$, integrating each shift
energy and applying \eqref{eq:kernel-identity} yields
\begin{equation}\label{eq:gamma-positive-boundary}
 t_\Gamma(u)=\sum_{n\ge0}d_\Gamma(k_n)|f_n|^2
 +\sum_{j\ge0}2(\coth(\alpha_jB)-1)
       |\langle f,\cosh(\alpha_jx)\rangle|^2.
\end{equation}
Every term in both series is nonnegative, so monotone convergence
justifies the identity without discarding any gamma energy.

Put $\epsilon_n=(-1)^n$. The second series has matrix
\begin{equation}\label{eq:boundary-matrix}
 \mathcal B_{nm}=b_nb_m\epsilon_n\epsilon_m
 \sum_{j\ge0}\frac{\alpha_j^2(1-e^{-2\alpha_jB})}
 { (\alpha_j^2+k_n^2)(\alpha_j^2+k_m^2)}.
\end{equation}
The pole equals $4|\langle f,\cosh(x/2)\rangle|^2$; its matrix is
$4cc^{\mathsf T}$, where
\begin{equation}\label{eq:pole-vector}
 c_n=\frac{b_n\epsilon_n\sinh(B/2)}{2(1/4+k_n^2)}.
\end{equation}
Thus the exact matrix of the signed form is
\begin{equation}\label{eq:Q-neumann}
 Q_{nm}=D_n\delta_{nm}+\mathcal B_{nm}+4c_nc_m-P_{nm}.
\end{equation}

These formulas can first be used on smooth compact tests. The finitely
many low cosine modes, zero extended and reflected evenly, belong to
$\mathcal D_B$: their translation increments have squared norm $O(s)$
near zero, so their gamma energy is finite. For a direct approximation,
cut off each endpoint over length $\varepsilon$ and smooth. The difference
has squared norm $O(\varepsilon)$, uniformly bounded variation and
uniformly bounded amplitude; hence its squared translation increment is
$O(\min(s,\varepsilon))$. Integration against $\rho$ is
$O(\varepsilon(1+|\log\varepsilon|))$. This proves convergence in the
minimum form norm. Subtracting the finite low-mode projection is therefore
legitimate. For each fixed low mode the cross coefficients below are
$O(n^{-1})$, so they represent an $L^2$ vector. More directly, on the
open interval the gamma operator is the integral of
$\rho(s)[2u(x)-u(x+s)-u(x-s)]$. For a zero-extended cosine, splitting
this integral at the distance to the boundary bounds its absolute value
by a constant times $1+|\log(B-|x|)|$. This is square integrable. Its pairing with any high
form-domain vector is consequently the stated convergent coefficient
pairing. All inequalities extend from the smooth core by form convergence.

\subsection{The active prime and the exact coupling}
Fix $B=1/2$ for the remainder of the proof. Then $k_n=2\pi n$,
$b_0=\sqrt2$, and $b_n=2$ for $n\ge1$. Since
$1/2<\log2<1<\log3$, the sole active prime power is $2$.
Put $s=\log2$, $a=(\log2)/\sqrt2$, and $x_0=s-B$.
On the positive half interval its operator is exactly
\[
 (Pf)(x)=a\,\mathbf1_{[x_0,B]}(x)f(s-x),\qquad \norm P=a.
\]
In particular, its signs have not been replaced by a positive form.
Its matrix is
\begin{equation}\label{eq:prime-cosine}
 P_{nm}=ab_nb_m\int_{x_0}^B
           \cos(k_nx)\cos(k_m(x-s))\dd x.
\end{equation}
For all $n,m$, put $\ell=2B-s$ and $\operatorname{sinc}t=\sin(t)/t$,
with $\operatorname{sinc}0=1$. The integral equals
\[
 \frac\ell2\left[
 \operatorname{sinc}\!\left(\frac{(k_n-k_m)\ell}2\right)
       \cos\!\left(\frac{(k_n+k_m)s}2\right)
 +\operatorname{sinc}\!\left(\frac{(k_n+k_m)\ell}2\right)
       \cos\!\left(\frac{(k_n-k_m)s}2\right)\right].
\]
For $n>m$, the equivalent endpoint expression is
\begin{equation}\label{eq:prime-endpoints}
 -\frac{k_n\epsilon_m\sin(k_nx_0)
          +k_m\epsilon_n\sin(k_m(B-s))}{k_n^2-k_m^2}.
\end{equation}
Both expressions follow by integrating the cosine product.

Take $M=16$. Split $f=z+h$ into modes $0\le m<M$ and $n\ge M$.
Let $A=(Q_{nm})_{n,m<M}$ and let $R$ have rows
$r_n=(Q_{nm})_{m<M}$ for $n\ge M$.
The full high block $T$ satisfies, as quadratic forms,
\begin{equation}\label{eq:high-diagonal}
 T\succeq K:=\operatorname{diag}_{n\ge M}K_n,
 \qquad K_n=D_n-a.
\end{equation}
This follows from the nonnegative boundary and pole forms and
$P\preceq aI$. The function $d_\Gamma(k)$ increases for $k\ge0$,
as is immediate from its series. The interval calculation verifies
$K_M>2$, hence $K\succeq2I$.
Consequently
\begin{equation}\label{eq:schur-frequency}
 \begin{split}
 q_{1/2}(u)&\ge
 z^{\mathsf T}(A-R^*K^{-1}R)z+
 \norm{K^{1/2}h+K^{-1/2}Rz}_{\ell^2}^2,\\
 R^*K^{-1}R&=\sum_{n\ge M}\frac{r_n^{\mathsf T}r_n}{K_n}.
 \end{split}
\end{equation}
The actual Schur complement is $A-R^*T^{-1}R$ and is at least
$A-R^*K^{-1}R$. The inverse cost retains its frequency dependence.
No claim is made to have evaluated $T^{-1}$ exactly.

\subsection{A rank-one bound for the entire omitted mode tail}
Set $N=8192$ and $v_m=b_m\epsilon_m$ for $m<M$. Define
\[
 U(k)=\sum_{j\ge0}\frac1{\alpha_j^2+k^2},\qquad U_n=U(k_n).
\]
For $n\ne m$ the nonexponential part of \eqref{eq:boundary-matrix} is
\[
 \sum_j\frac{\alpha_j^2}{(\alpha_j^2+k_n^2)(\alpha_j^2+k_m^2)}
 =U_n+\frac{k_m^2(U_n-U_m)}{k_n^2-k_m^2}.
\]
The elementary decreasing-sum bound gives, for $k>0$,
\begin{equation}\label{eq:Ubound}
 U(k)\le\frac1{k^2+1/4}+\frac\pi{4k}.
\end{equation}
Indeed, keep the first summand and bound the rest by one half of the
integral of $(t^2+k^2)^{-1}$ on $t\ge1/2$.

Let $E_0=e^{-B}/(1-e^{-4B})$, and define
\begin{align*}
 H_0&=0,\qquad H_m=1+\pi k_m/4\quad(m\ge1),\\
 C_m&=b_m\left[
 \frac{H_m}{1-(k_m/k_N)^2}+E_0+
 a\frac{k_m+k_m^2/k_N}{1-(k_m/k_N)^2}\right]
       +2\sinh(B/2)|c_m|.
\end{align*}
For every $n\ge N$ the exact row admits
\begin{equation}\label{eq:tail-row}
 r_n=\alpha_n v^{\mathsf T}+e_n^{\mathsf T},\qquad
 \alpha_n=b_n\left[\epsilon_nU_n+
                          \frac{a\sin(k_nx_0)}{k_n}\right],
 \qquad |e_{nm}|\le\frac{b_n C_m}{k_n^2}.
\end{equation}
Here $e_n$ denotes a remainder vector, not a cosine basis function.
To check the bound, use $0\le U_n\le U_m$ for $n>m$;
$k_m^2U_m\le H_m$; and
\[
 \sum_j\frac{\alpha_j^2e^{-2\alpha_jB}}
 { (\alpha_j^2+k_n^2)(\alpha_j^2+k_m^2)}\le\frac{E_0}{k_n^2}.
\]
Subtracting the leading term in \eqref{eq:prime-endpoints} leaves at most
$a b_nb_m(k_m+k_m^2/k_N)/(k_n^2(1-(k_m/k_N)^2))$.
Equation \eqref{eq:pole-vector} gives the last term in $C_m$.
These account for the complete cross row, including its signs.

The verifier checks $K_N>1$; monotonicity gives $K_n>1$ for $n\ge N$.
By $(x+y)(x+y)^{\mathsf T}\preceq2xx^{\mathsf T}+2yy^{\mathsf T}$,
\eqref{eq:Ubound}, and the integral bounds for $\sum n^{-2}$ and
$\sum n^{-4}$, we obtain
\begin{equation}\label{eq:tail-bound}
 \sum_{n\ge N}\frac{r_n^{\mathsf T}r_n}{K_n}
 \preceq L_N vv^{\mathsf T}+E_N I,
\end{equation}
where the following explicit, deliberately conservative constants suffice:
\begin{equation}\label{eq:tail-constants}
 L_N=\frac{8(\pi/4+a+1/k_N)^2}{(2\pi)^2(N-1)},\qquad
 E_N=\frac{8\sum_{m<M}C_m^2}{3(2\pi)^4(N-1)^3}.
\end{equation}
Retaining the rank-one direction $v$ is useful: replacing this term by
its scalar operator norm would lose substantially more information.

\subsection{Finite interval certificate and the final norm conversion}
Define the exact finite matrix
\begin{equation}\label{eq:Cbound}
 \mathcal C=\sum_{n=M}^{N-1}\frac{r_n^{\mathsf T}r_n}{K_n}
                  +L_Nvv^{\mathsf T}+E_NI.
\end{equation}
The preceding analytic estimates prove
$R^*K^{-1}R\preceq\mathcal C$.
The finite entries are evaluated using
\[
 U(0)=\psi_1(1/4)/4,\qquad
 U(k)=\frac{\Im\psi(1/4+ik/2)}{2k}\quad(k>0).
\]
\[
 \sum_j\frac{\alpha_j^2}{(\alpha_j^2+k^2)^2}
 =\frac{\Im\psi(1/4+ik/2)}{4k}
       +\frac{\Re\psi_1(1/4+ik/2)}8\quad(k>0).
\]
At $k=0$ the last sum is $U(0)$. Here $\psi_1=\psi'$.
These expressions follow from the same absolutely convergent series
and its derivative. Only the exponentially weighted subtraction in
\eqref{eq:boundary-matrix} is truncated. Keeping $0\le j<J$, $J=24$,
leaves an entrywise inner-sum error bounded by
\begin{equation}\label{eq:exp-tail}
 \frac{e^{-2\alpha_JB}}{\alpha_J^2(1-e^{-4B})}.
\end{equation}
The factors $b_nb_m\epsilon_n\epsilon_m$ are applied after enclosing
that error.

The supplied script uses 160-bit real and complex ball arithmetic
from \texttt{python-flint} 0.9.0 \cite{pythonflint}. All mathematical
inputs are integers or ratios of integers. It encloses every entry of
$A$ and $\mathcal C$, including \eqref{eq:exp-tail}, and applies interval
$LDL^{\mathsf T}$ elimination to
\[
 A-\mathcal C-\mu I,\qquad \mu=8/10^7.
\]
Every pivot has a strictly positive lower endpoint; no approximate
eigenvalue is used as a certificate. The recorded run verifies
\begin{equation}\label{eq:certified-tests}
 K_M>2,\qquad K_N>1,\qquad
 \operatorname{tr}\mathcal C<1/10,\qquad A-\mathcal C\succeq\mu I.
\end{equation}
Some outward-rounded enclosures from the run are:
\begin{center}
\begin{tabular}{@{}ll@{}}
\toprule
Quantity&Enclosure\\\midrule
$K_M$&$(2.2824555276,\ 2.2824555278)$\\
$L_N$&$(4.0251866750,\ 4.0251866752)\,10^{-5}$\\
$E_N$&$(1.0342902112,\ 1.0342902113)\,10^{-9}$\\
$\operatorname{tr}\mathcal C$&$(0.09725200405,\ 0.09725200407)$\\
Smallest interval-$LDL$ pivot&$(0.0015928668,\ 0.0015928670)$\\
\bottomrule
\end{tabular}
\end{center}
An $LDL$ pivot is not an eigenvalue. The positive pivots establish
positive definiteness of the matrix to which elimination was applied.
Full enclosures and the verifier hash are in
\texttt{half\_support\_certificate.json}.

It remains to convert the low-mode margin into an ordinary-$L^2$ bound
on all inputs. Since $K\succeq2I$,
\[
 \norm{K^{-1}R}^2\le\tfrac12\norm{R^*K^{-1}R}
 \le\tfrac12\operatorname{tr}\mathcal C<1/20.
\]
Set $t=\norm{K^{1/2}h+K^{-1/2}Rz}^2$. Then
\[
 q_{1/2}(u)\ge\mu|z|^2+t,\qquad
 \norm u_2^2=|z|^2+\norm h^2\le\tfrac{11}{10}|z|^2+t.
\]
Because $(7/10^7)(11/10)<\mu$ and $7/10^7<1$, this proves
\eqref{eq:halfwin}. For $B<1/2$, the defining core and its closure embed
in $\mathcal D_{1/2}$ with the identical zero-extended form and norm,
so the same lower bound follows. This completes the computer-assisted
proof of Theorem~\ref{thm:half}.

For any continuous positive weight $W$ on $[-1/2,1/2]$, put
$w_* =\min W>0$. A valid conversion, if needed, is
\[
 q_{1/2}(u)\ge(7/10^7)w_* N_X(u),
 \qquad N_X(u)=\int |u|^2/W.
\]
The coefficient $7/10^7$ by itself belongs to ordinary $L^2$.
It cannot be added to bounds for other operators or metrics.

\section{Why a vanishing all-domain defect is not a weaker target}
For a fixed positive weight $W$, define the extended nonnegative number
\[
 d_B^*=\max\left\{0,\sup_{0\ne u\in\mathcal D_B}
                    \frac{-q_B(u)}{N_X(u)}\right\}.
\]
The domains are nested and the zero-extended form and norm agree on
overlaps. Hence $d_{B_1}^*\le d_{B_2}^*$ for $B_1\le B_2$.
If $B_j\to\infty$ and full-domain bounds
$q_{B_j}\ge-d_jN_X$ hold with $d_j\to0$, fix any smaller-domain input
and send $j\to\infty$. Its form value is nonnegative. Conversely,
positivity for every compact domain permits $d_j=0$.
Thus the vanishing-defect family is equivalent to positivity on every
compact local domain. Theorem~\ref{thm:half} proves $d_B^*=0$ for
$B\le1/2$ only. It does not establish that family at unbounded radii.

Beyond $B=1/2$, the same gamma representation remains available but
the active prime shifts and their coupling rows change. As the radius
crosses $(\log3)/2,(\log4)/2,\ldots$, further prime powers enter.
The next substantive requirement is a new sufficient Schur certificate
including all these rows, or an analytic bound uniform in the radius.
Even a certificate at one larger radius would still be a compact-support
result. A localization error tending to zero supplies no automatic sign
for an unresolved interior Schur complement.

\section{What this proves, and what remains open}
\begin{enumerate}
\item Theorem~\ref{thm:half} supplies a full-domain local positive lower
bound for $B\le1/2$ by an analytic reduction plus a reproducible interval
calculation. It is not a historical novelty claim. The software, analytic
derivation and implementation remain subject to independent audit.
\item The local nonlinear and quadratic inequalities have direct analytic proofs. They require no RH, numerical spectral certificate, prime number theorem or evaluation of a theta source. They apply to the complete compact-support signed form.
\item Equation \eqref{eq:global} additionally assumes \eqref{eq:localization-hypothesis}. The ordinary-$L^2$ coefficient $A_{B,\delta}$ is not automatically a spectral lower bound in a differently weighted Hilbert space.
\item Moment vanishing forces a primitive to vanish at both endpoints of every cell. Dirichlet Poincare then suppresses the low-frequency resolvent norm; the logarithmic gamma symbol converts that suppression into coercivity. This is the structural reason for domination on the stated class.
\item At arbitrary radii, general inputs can have large $P_\delta u$. A positive infinite-dimensional moment kernel does not prove positivity of the full form. The finite quotient still requires its exact Schur complement, including coupling with that kernel. It is not enough to check the raw moment-span restriction.
\item Refining cells does not force the sufficient bound to succeed on a fixed general input. Typically $P_\delta u\to u$ under refinement; the quadratic penalty may become very large. The nonlinear bound avoids that artificial divergence but still supplies no global positivity. Increasing $R$ also changes $\beta_B$, the partition, and the finite quotient.
\item There is no claimed uniform lower bound on interior mass for all inputs. No positive lower bound for $t_\Gamma(\eta_Ru)$ in the $N_X$ norm has been established here. Thus a local margin larger than $\varepsilon_R$ alone is not a global proof.
\end{enumerate}

For a global conclusion through localization, one possible route is to prove \eqref{eq:localization-hypothesis} with
$\varepsilon_R\to0$ along an unbounded sequence, and full-domain local
bounds $q_B(v)\ge-\delta_BN_X(v)$ with $\delta_B\to0$ along the
corresponding radii $B=R+1$. Since $N_X(\chi_Ru)\le N_X(u)$ and
$t_\Gamma(\eta_Ru)\ge0$, these would give
\[
 q_\infty(u)\ge-(\delta_{R+1}+\varepsilon_R)N_X(u),
\]
and then $q_\infty(u)\ge0$ by taking the limit for each fixed test.
The vanishing all-domain local defects would already imply positivity
on every compact support without this localization step, by domain
nesting. Thus localization is not a shortcut around that requirement.
Extension to a chosen minimum closed domain requires its core-density
and form-convergence justification. The necessary family of local
bounds is not proved here. No identification with a differently
normalized spectral operator is asserted.

\section{Checks and scope}
The companion exact-check script verifies the projection Gram determinant
identity, the coefficients in the bounded-resolvent argument, the Fourier
denominator cancellation, the exponential-moment annihilation identities,
and the finite Schur square identity. These are finite algebra checks,
not a proof checker for the analytic argument. The additional interval verifier certifies precisely the finite
inequalities in \eqref{eq:certified-tests}, using the analytic remainder
bounds proved in Section~\ref{sec:half}. No binary64 eigenvalue is used
as evidence of a sign. The PDF contains the TeX source, both check
scripts, and their outputs as attachments for inspection.
The references below supply standard inequality and closed-form
background; they do not establish the application-specific domination
theorem or the localization hypothesis.

\paragraph{AI assistance.}
Artificial intelligence assistance was used to draft and check the
mathematical exposition and prepare the manuscript, finite algebra checks and interval verifier. This statement does not assert external peer review or
proof-assistant verification.

\begin{thebibliography}{9}
\bibitem{HLP}
G.~H. Hardy, J.~E. Littlewood, and G. P\'olya,
\emph{Inequalities}, 2nd ed., Cambridge University Press, Cambridge,
1952; Cambridge Mathematical Library reprint, 1988.
ISBN 978-0-521-35880-4.
\href{https://assets.cambridge.org/97805213/58804/frontmatter/9780521358804_frontmatter.pdf}{Publisher's front matter};
\href{https://assets.cambridge.org/97805213/58804/toc/9780521358804_toc.pdf}{table of contents}.

\bibitem{Kato}
T. Kato, \emph{Perturbation Theory for Linear Operators}, 2nd ed.,
Classics in Mathematics, Springer-Verlag, Berlin--Heidelberg, 1995.
\href{https://doi.org/10.1007/978-3-642-66282-9}{doi:10.1007/978-3-642-66282-9}.

\bibitem{DLMF}
NIST Digital Library of Mathematical Functions,
Chapter 5, Gamma Function, especially \S5.7(ii),
\href{https://dlmf.nist.gov/5.7}{https://dlmf.nist.gov/5.7}.

\bibitem{pythonflint}
The python-flint contributors,
\emph{python-flint 0.9.0 documentation},
real and complex ball arithmetic and special functions,
\href{https://python-flint.readthedocs.io/en/stable/arb.html}{arb documentation};
\href{https://python-flint.readthedocs.io/en/stable/acb.html}{acb documentation}.
\end{thebibliography}
\end{document}
